\chapter{Task B. Patient-Level Endpoint Prediction}
\label{ch:taskB-method}
\label{ch:taskB-eval}

\section{Problem Formulation}
\label{sec:taskB-problem}

Task B is formulated as multi-label node classification on a heterogeneous,
per-patient graph, in order to predict clinical outcomes.
Instead of building one global graph and treating patients as additional nodes,
every patient $p$ is represented as a copy of the same fixed SPG~v3, with node
feature vectors derived from that patient's covariates.

Formally, patient $p$ is represented by the attributed graph
$G_p = (V, E, x_p)$, where $x_p$ denotes the patient-specific node features.
The topology $(V, E)$ is shared across the cohort.
The node set $V$ comprises six semantic categories:
\begin{itemize}
    \item \textit{patient\_context} --- baseline characteristics such as age, sex, or alcohol use;
    \item \textit{drug\_exposure} --- medications taken, such as anticoagulants, diuretics, or opioids;
    \item \textit{mechanism} --- biological or pharmacodynamic processes such as slowed drug breakdown, sodium--water imbalance, or poor kidney blood flow;
    \item \textit{adr\_or\_intermediate\_state} --- visible medical states such as a rise in creatinine, ECG-measured QT prolongation, or bleeding;
    \item \textit{observation\_or\_selection} --- testing, documentation, and care-selection variables such as laboratory monitoring or recorded clinical checks;
    \item \textit{clinical\_endpoint} --- clinical outcomes being predicted, described in \cref{subsec:taskB-prediction-targets}.
\end{itemize}
These node types encode a causal pathway towards a clinical outcome:
signals propagate along the edges of SPG~v3 from context and drug exposure,
through mechanisms and visible medical states, to the endpoints.
In general, heart failure (\textit{patient\_context}) together with an NSAID
(\textit{drug\_exposure}) may reduce kidney blood flow (\textit{mechanism}),
become visible as a rise in creatinine
(\textit{adr\_or\_intermediate\_state}), and culminate in AKI
(\textit{clinical\_endpoint}); laboratory monitoring of creatinine is an
\textit{observation\_or\_selection} node on the same pathway.
Which of those nodes are actually on for a given patient is encoded in $x_p$.
\Cref{fig:patient_dag} shows one such realisation ending in AKI: filled nodes
are this patient's active context, drugs and mechanisms, whereas outlined nodes
(for example NSAID, heart failure, or a creatinine rise) remain inactive
mediators on the shared schema.
The aim of Task~B is to estimate endpoint
risk from the entire patient-specific graph as described above.

\begin{figure}[H]
    \centering
    \includegraphics[width=\textwidth]{patient_dag_101.pdf}
    \caption{Exemplary per-patient causal DAG for one synthetic patient (id~101).
    The full SPG~v3 is shared across the cohort; this figure keeps only the
    AKI-relevant fragment so that the left-to-right causal layers remain readable.
    Filled nodes are active for this patient; outlined nodes are inactive mediators
    on the same paths.}
    \label{fig:patient_dag}
\end{figure}

Assuming that \(\mathcal{V}_{\text{target}} \subset V\) is the set of
endpoint nodes used as prediction targets,
for each patient $p$ and target node $v \in \mathcal{V}_{\text{target}}$
the binary label $y_{p,v}$ indicates whether the endpoint occurs.
A shared heterogeneous GNN $f_\theta$ produces a separate probability for each
endpoint from the same encoded graph; $v$ is a concrete clinical endpoint such
as QT arrhythmia, Falls, Serotonin syndrome, or AKI:
\begin{equation}
\begin{gathered}
\hat{y}_{p,v} = \sigma\big(f_\theta(G_p)_v\big),
\quad v \in \mathcal{V}_{\text{target}}, \\[0.4em]
{\scriptstyle
G_p=(V,E,x_p),\ 
f_\theta\text{ --- heterogeneous GNN shared across patients},\ 
v\text{ --- clinical endpoint node}
}
\end{gathered}
\end{equation} 

\section{Patient-Level Graph Representation}
\label{subsec:taskb_data}

The input graph for the node classification task is built from: 

\begin{enumerate}[label=\alph*)]
    \item \textbf{nodes}
    \item \textbf{edges}
    \item \textbf{per-scenario patient sample table} 
\end{enumerate}

These components are assembled into a heterogeneous per-patient graph representation
 whose elements are summarised in \cref{tab:heterodata_components}.

\begin{table}[H]
    \centering
    \caption{Components of the per-patient heterogeneous graph representation.}
    \label{tab:heterodata_components}
    \begin{tabularx}{\textwidth}{@{}l X@{}}
        \toprule
        \textbf{Component} & \textbf{Description} \\
        \midrule
        Node types &
        Six semantic categories of nodes shared across all patient graphs: \texttt{patient\_context}, \texttt{drug\_exposure}, \texttt{mechanism}, \texttt{adr\_or\_intermediate\_state}, \texttt{observation\_or\_selection}, and \texttt{clinical\_endpoint}. \\
        \addlinespace
        Relations &
        Keyed by \texttt{(src\_type, "to", dst\_type)};
        \texttt{edge\_type} is encoded as a one-hot block inside
        \texttt{edge\_attr}. This keeps the number of distinct convolutional
        relation modules in \texttt{HeteroConv} small while still letting
        edge-type information reach the message function. \\
        \addlinespace
        Edge attributes (\texttt{edge\_attr}) &
        Fixed-length vector per edge (see \cref{tab:edge_attr_layout}): signed
        effect \texttt{effect\_size}$\times$\texttt{effect\_sign}, audit scalars
        \texttt{activation\_frequency}, \texttt{mechanistic\_confidence},
        \texttt{evidence\_weight}, and a one-hot \texttt{edge\_type} block. \\
        \addlinespace
        Node features (\texttt{x}) &
        The patient's observed values for a node \\
        \bottomrule
    \end{tabularx}
\end{table}

Patient observations are encoded as node features.
The AKI fragment in \cref{fig:patient_dag} is the running example of
Section~\ref{sec:taskB-problem}: filled nodes follow this patient's $x_p$,
outlined nodes are inactive mediators on the same directed paths.

Two structural exclusion mechanisms are applied before any feature is computed: 
\begin{enumerate}[label=\alph*)]
    \item nodes flagged \texttt{exclude\_from\_observed\_model} or \texttt{is\_latent} (hidden confounders); 
    \item nodes that are structural descendants of any endpoint in the DAG (e.g. documentation nodes such as 
    \texttt{hospital\_contact}, \texttt{fall\_reported}, which are caused by the endpoint rather than causal parents of it)
\end{enumerate}

 All normalization statistics (mean/std for continuous node values) 
 and all statistical evidence features (\cref{subsec:hcr_application}) are computed exclusively on the training split of patients, to 
 prevent both node-level and patient-level information leakage across splits.

\subsection{Prediction targets}
\label{subsec:taskB-prediction-targets}
Although the knowledge graph defines 13 
clinical outcomes nodes, the model is trained and evaluated on a
subset of 11 clinical endpoints, excluding
%\texttt{AKI}, \texttt{DILI}, \texttt{Depression}, \texttt{Delirium},
%\texttt{GI\_bleeding}, \texttt{Hyponatremia}, \texttt{Hyperkalemia},
%\texttt{QT\_arrhythmia}, \texttt{Serotonin\_syndrome},
%\texttt{Rhabdomyolysis}, and \texttt{Lactic\_acidosis}.
%The remaining two endpoints---
\texttt{Hospitalization} and
\texttt{Falls}. They are retained in the per-patient graph topology and in
patient labels, but are excluded from the multi-target prediction head.

This choice follows both the graph structure and the clinical scope of
the task. \texttt{Hospitalization} was excluded because the DAG models it
as a downstream consequence of other endpoints
(\texttt{endpoint\_to\_hospitalization}), which makes it a poor fit for
parallel multi-target prediction alongside those same outcomes. 
\texttt{Falls}, by contrast, is a geriatric outcome that has different 
causal pathways than other endpoints that originate from exposures
and mechanistic intermediates in the pharmacotherapy graph. On these grounds, Falls and Hospitalization are treated as downstream, multifactorial events outside
that primary prediction scope.

For reporting, all eleven targets are trained, but macro-averaged metrics
(AUROC, AP, \% of oracle AUROC; \cref{eq:pct-ceiling}) are computed on eight endpoints only.
\texttt{Serotonin\_syndrome}, \texttt{Rhabdomyolysis}, and
\texttt{Lactic\_acidosis} are omitted from the aggregate because their
extreme rare prevalence makes
split-level AUROC and AP unstable.
The resulting reporting hierarchy is therefore 13 graph endpoints $\rightarrow$ 11 trained targets $\rightarrow$ 8 endpoints included in macro metrics.

\section{Model Architecture}
\label{sec:taskB-architecture}

The core architecture is a stack of \texttt{HeteroConv} layers (\texttt{HeteroConv} layers (a wrapper applying a homogeneous message-passing operator per 
relation type; \cref{subsec:related-gnn}) 
applied to the
SGP~v3. 
Each layer instantiates a separate message-passing operator per relation type; 
the same backbone weights are shared across all patient graphs
in the cohort. Three operators 
were benchmarked against each other: GraphSAGE, R-GCN, and
TransformerConv.
All three are applied to SGP~v3 via \texttt{HeteroConv}.
Each directed edge in the graph carries a fixed \texttt{edge\_attr} vector 
assembled at
graph-build phase from the audited edge table (\cref{tab:edge_attr_layout}).
In model design relations are grouped by node-type pairs \texttt{(src\_type, to, dst\_type)}
and the semantic \texttt{edge\_type} is
retained as a one-hot suffix in \texttt{edge\_attr}. 
This vector is consumed differently depending on the message-passing
operator (\cref{tab:edge_attr_by_conv}).

\begin{table}[H]
  \centering
  \small
  \caption{Layout of the Task~B \texttt{edge\_attr} vector.}
  \label{tab:edge_attr_layout}
  \begin{tabularx}{\textwidth}{@{}l l X@{}}
    \toprule
    \textbf{Block} & \textbf{Columns} & \textbf{Source} \\
    \midrule
    Signed effect &
    1 &
    \texttt{effect\_size} $\times$ \texttt{effect\_sign} \\
    Audit scalars &
    3 &
    \texttt{activation\_frequency}, \texttt{mechanistic\_confidence},
    \texttt{evidence\_weight} \\
    Edge-type one-hot &
    $|\texttt{edge\_type}|$ &
    semantic edge category from the audited DAG \\
    \bottomrule
  \end{tabularx}
\end{table}

\begin{table}[H]
  \centering
  \small
  \caption{Use of \texttt{edge\_attr} by convolution backbone in Task~B.}
  \label{tab:edge_attr_by_conv}
  \begin{tabularx}{\textwidth}{@{}l X@{}}
    \toprule
    \textbf{Backbone} & \textbf{Edge features used in message passing} \\
    \midrule
    GraphSAGE &
    none (topology only; \texttt{edge\_attr} stored but not passed to
    \texttt{SAGEConv}) \\
    TransformerConv &
    full \texttt{edge\_attr} via \texttt{edge\_dim} \\
    R-GCN &
    relation identity via separate modules per edge type, not via the
    Task~B \texttt{edge\_attr} vector \\
    \bottomrule
  \end{tabularx}
\end{table}

A design issue handled during the experiments is the duplication of the self-transform 
across relations. 
Several PyG message-passing layers (SAGEConv, TransformerConv) internally add their 
own linear self-loop term,
 wrapped per-relation inside \texttt{HeteroConv}. This contribution of the node’s 
 own representation is duplicated once 
 per incoming relation type. The implementation applied removes this internal duplication 
(internal self-loop weights disabled for SAGE/Transformer) and replaces it with a single, 
explicit, 
controlled self-transform per layer per node type. 
% For GAT-family layers, where the self-dependency is baked into the attention query and cannot be disabled, the explicit self-transform 
% acts as an additional controlled term on top of the irreducible built-in dependency. GraphConv/GCN, which has no 
% \texttt{root\_weight} parameter at all, is deliberately excluded from this mechanism, since its built-in self-transform 
% cannot be removed without reimplementing the layer.
%node identity 

The architecture begins with the linear layer that maps each node’s raw features into a hidden vector of fixed width (the input projection). 
Several clinical endpoints in SPG~v3---serotonin syndrome, rhabdomyolysis, and lactic acidosis---share identical 
audited attributes and similar neighbourhoods, so a message-passing GNN cannot recognise them apart. 
This limitation is the one-dimensional Weisfeiler--Lehman colouring: nodes with equal features and equivalent neighbourhoods 
receive the same representation~\cite{wojcik2026grafowe}. To break the symmetry, a learnable node-identity 
embedding $x'$ is added after the input projection and before the first message-passing layer:
% To break this symmetry, a learnable embedding table is maintained per node type,
% indexed by each node's fixed local index in the shared DAG, and added
% to the projected input features before the first message-passing layer:
\begin{equation}
\begin{gathered}
x' = Wx + e_{\mathrm{type}}[\mathrm{index}], \\[0.3em]
{\scriptstyle
x\text{ --- raw node features},\ 
W\text{ --- input projection},\ 
e_{\mathrm{type}}[\mathrm{index}]\text{ --- identity vector for that node in the shared DAG}
}
\end{gathered}
\label{eq:node-identity}
\end{equation}
% Equation~\eqref{eq:node-identity} is applied after the input projection and before message passing; the endpoint head
%  reads the resulting node representations.
For each node type, a table of vectors is indexed by the node's fixed local index
$e_{\mathrm{type}}[\mathrm{index}]$ in SPG~v3 and
the looked-up vector is added to the projected features so that similar endpoints remain 
distinguishable even
when their static attributes coincide. The sequence is presented in \cref{fig:node-identity}.

\begin{figure}[H]
\centering
\resizebox{\textwidth}{!}{%
\begin{tikzpicture}[
  font=\small,
  box/.style={draw, rounded corners=2pt, align=center, inner sep=5pt, minimum height=1.35cm},
  arr/.style={-{Stealth}, thick}
]
\node[box, fill=black!6] (raw) {Raw features $x$\\[-1pt]{\scriptsize identical for three endpoints}};
\node[box, fill=blue!8, right=0.55cm of raw] (proj) {Input projection\\[-1pt]$Wx$};
\node[box, fill=orange!15, right=0.55cm of proj] (id) {Node-identity\\[-1pt]embedding $e[\mathrm{index}]$};
\node[box, fill=green!10, right=0.55cm of id] (mp) {Message-passing\\[-1pt]layers (\texttt{HeteroConv})};
\node[box, fill=black!6, right=0.55cm of mp] (head) {Endpoint\\[-1pt]head};

\draw[arr] (raw) -- (proj);
\draw[arr] (proj) -- (id);
\draw[arr] (id) -- (mp);
\draw[arr] (mp) -- (head);

\node[below=0.45cm of id, font=\scriptsize, align=center] (tab) {%
  lookup in shared DAG:\\[-1pt]
  serotonin syndrome $\mapsto e_7$,
  rhabdomyolysis $\mapsto e_8$,
  lactic acidosis $\mapsto e_9$};
\draw[arr, thin] (tab) -- (id);
\end{tikzpicture}%
}
\caption{Order of operations in the Task~B encoder.}
\label{fig:node-identity}
\end{figure}
%oversmoothing
\label{subsubsec:oversmoothing}

A separate issue at greater depth is oversmoothing. When several message-passing layers are stacked,
repeated aggregation can make node representations increasingly similar~\cite{rusch2023oversmoothing}.
Residual connections, Jumping Knowledge, DropEdge, and PairNorm were tested as countermeasures;
residual connections were the most effective and are used in the results reported below.

%These methods were evaluated both 
%individually and in combination with the HCR enrichment described in \cref{subsec:hcr_application}.

% \section{Statistical evidence branch}
% \label{subsec:hcr_application}

% In addition to the purely structural GNN pathway, approch presented in the thesis incorporates 
% a statistical evidence 
% side-channel with dependecies inspired by HCR described in \cref{subsec:hcr_def}. 
% This channel is motivated by the 
% observation that some parent-endpoint dependencies are well captured by simple, 
% well-understood pairwise 
% statistics computed directly from the training cohort, and that combining such statistics
%  with the deep GNN 
% representation captures the strong signal from direct parents and smoothed signal 
% from further layers.
% There were run couple of experiments with HCR configurations and GNNs. However, in general the approach was about
% computation of contingency-table statistics for each target endpoint and direct parents or pre-direct parents in 
% the audited DAG (restricted to nodes with inferred 
% binary, count, or continuous value types). The statistics computed was as followed per (parent, endpoint) 
% pair: the \(\varphi\)-coefficient, normalized mutual information, a tanh-scaled log odds 
% ratio, risk difference, and a tanh-scaled log lift, plus prevalence and support terms. There 
% are nine to forty-one 
% channels depending on configuration.
% %dopisac dlaczego u nas jest 41/42 a w przypadku task a 41
% %leave-one-out

% As these coefficients are computed with the patient's own endpoint label as part of the statistic, a naive 
% computation would leak the label into the feature. This was addressed with an 
% exact leave-one-out correction: every training patient's coefficient is computed from all other training patients
% only, using a closed-form re-weighting of the population statistic.
 
%  Continuous and count-valued parents are mapped to comparable pseudo-uniform scores via a train-only empirical 
%  CDF transform (with deterministic, patient-ID-seeded jitter for count variables), and Legendre polynomial bases 
%  are used to capture non-linear dependence on continuous parents.

%  Four integration modes were supported: 
%  \begin{enumerate}[label=\alph*)]
%     \item linear approach - a single scalar HCR score added directly to the GNN logit,
%     \item nonlinear encoder - parent-endpoint evidence vectors pooled and passed through a shared MLP
%     \item separate encoders per parent node type - every direct parent passed through a separated MLP
%     \item triple relations - relations not only focused on direct parents, but direct parents and grandparents. 
%  \end{enumerate}

%  The pooling in HCRPairEncoder masks out padding positions (endpoints with fewer parents than the batch's 
%  maximum) and returns a neutral zero vector for endpoints with no eligible parents, avoiding division by zero.

%jak jest wszystko polaczone dokladniej

\section{Statistical Evidence Branch}
\label{subsec:hcr_application}

In addition to the deep graph representation learned by the GNN, endpoint
prediction includes a statistical evidence branch: a wide side-channel that
supplies cohort-level pairwise dependence statistics.
The construction is \emph{HCR-inspired} (\cref{subsec:hcr_def}):
some parent--endpoint dependencies are well captured by interpretable
statistics computed on the training cohort, and combining them with the GNN
representation keeps a direct-parent signal that message passing can dilute.
% The exposition follows the conceptual progression of pairwise dependence,
% type-specific evidence encoding, higher-order triple-path evidence, and
% fusion with the GNN representation.
%Each mode is described in the subsections that follow.
%The HCR construction---a cohort coefficient multiplied by a
%patient-specific activation and summed over parents---is retained as one
%implementation mode (\emph{linear}).
% Figure~\ref{fig:taskB-hcr-arch} summarises the dual-path layout
% : a heterogeneous GNN producing endpoint embeddings, and an 
% HCR-inspired evidence branch that can be concatenated into the 
% classification head or added at the logit level.

\begin{figure}[H]
\centering
\includegraphics[width=\textwidth]{architecture_slide_multilabel_hcr.png}
\caption{Task~B architecture: heterogeneous GNN path with a shared endpoint head, 
HCR-inspired evidence tensors with pair encoders, and optional concatenation or 
logit-level fusion.}
\label{fig:taskB-hcr-arch}
\end{figure}

% The remaining configurations extend this idea with additional
% contingency-table statistics, non-binary parent handling, higher-order path
% moments, and learned fusion with the GNN head.
Figure~\ref{fig:taskB-hcr-arch} summarises the dual-path layout.
The GNN produces an endpoint embedding \(z_{\mathrm{graph}}\).
In the nonlinear and typed configurations, a pair encoder (a small MLP)
maps each parent--endpoint evidence vector to an embedding and
mean-pools over parents.
There may be a single encoder for all parents, or one encoder per parent
type (drug, mechanism, ADR, patient context) --
\(e_{\mathrm{drug}}\), \(e_{\mathrm{mech}}\), \(e_{\mathrm{adr}}\) and
\(e_{\mathrm{ctx}}\) presented at the bottom of the figure.
The resulting embeddings are either concatenated with \(z_{\mathrm{graph}}\)
as input to the classification head, or reduced to a scalar and added to the
GNN logit.
The different fusion modes tested are specified below (\cref{tab:fusion_modes}) 
and subsections below represents detailed info about configurations.


\subsection{Scope of variables and parents}
For each target endpoint, evidence is computed over its \emph{direct parents}
in the patient graph (subject to the structural exclusions described in
\cref{subsec:taskb_data}).
Pairwise statistics relate one parent to one endpoint at a single hop.
Higher-order evidence over \emph{grandparent--parent--endpoint} paths
(\(G \rightarrow P \rightarrow Y\), where \(G\), \(P\), and \(Y\) denote the
grandparent, parent, and endpoint nodes on SGP~v3).
is computed separately in the \emph{triple} extension.
It is stored in its own evidence tensor (13 channels per triple slot) and fused
through a dedicated encoder block.
Several evidence dimensionalities were evaluated, including binary-only
parents (compact 10-channel evidence) and all parent value types (binary,
count, and continuous; 42-channel evidence).

\subsection{Pairwise HCR core coefficients}
For a binary parent \(P\) and a binary endpoint \(Y\), the primary association
coefficient is the \(\varphi\)-coefficient
%(equivalently denoted \(a_{11}\) in the implementation)
derived from the training-cohort contingency table.
For each patient \(p\), the parent activation is the standardized binary score
\(\phi_1(R)\), where \(R\) denotes the parent.
The HCR patient evidence for that parent--endpoint pair is the product
\(s_{p,rv} = a_{11,rv}\,\phi_1(R_p)\) (channel \texttt{s\_evidence}).
In \emph{linear} mode the endpoint-wide score is the sum over eligible direct
parents,
\[
  s_{p,v} = \sum_{r \in \mathrm{Pa}(v)} a_{11,rv}\,\phi_1(R_p),
\]
and a single learned scalar per endpoint scales this score before it is
added to the GNN logit. 
%This is the only configuration that implements HCR
%in the strict sense; all other modes treat \(a_{11}\) and \(\phi_1\) as
%channels inside a richer evidence vector.

\subsection{Extended pairwise evidence channels}
Beyond the HCR core, the same contingency table yields supplementary
association measures that
provide complementary views of parent--endpoint dependence for a learned
encoder: normalized mutual information, a tanh-scaled log odds ratio, risk
difference, tanh-scaled log lift, parent and endpoint prevalence, and
support (the fraction of training patients with both variables observed).
These are assembled into a fixed-order evidence vector per
(patient, parent, endpoint) slot. Two pairwise dimensionalities were used:
\begin{itemize}
  \item \textbf{10 channels}: compact
        binary-parent representation with explicit \texttt{s\_evidence};
  \item \textbf{42 channels}:
        enriched representation including matrix-energy summaries, marginal
        entropies, type indicators, and the patient activation channel.
\end{itemize}
\subsection{Type-specific encoders and endpoint-specific fusion}
In \emph{nonlinear} and \emph{triple} modes, a shared pairwise evidence encoder 
(PairEncoder in \cref{fig:taskB-hcr-arch})
processes each evidence slot with a two-layer MLP
(\texttt{Linear}$\rightarrow$\texttt{ReLU}$\rightarrow$\texttt{Linear}),
mean-pools over parent or triple slots using a padding mask, and collapses
the pooled embedding to a scalar logit contribution. ReLU activation
is used as in the graph encoder. The MLP widths (hidden 32, output 8) were 
selected after preliminary runs with a narrower encoder (hidden 8, output 4) 
performed worse; they were then held fixed (\cref{tab:hparams-common}). 
In \emph{nonlinear} mode only, that contribution is further scaled by a
learned per-endpoint weight.

In \emph{typed\_concat} (and \emph{typed+triple}), the architecture differs
in two respects.
First, evidence is partitioned by \emph{parent node type}: each of the four
parent categories
(\texttt{drug\_exposure}, \texttt{mechanism},
\texttt{adr\_or\_intermediate\_state}, \texttt{patient\_context}) is encoded
by its own pairwise evidence encoder applied to a separate evidence
tensor (type-specific evidence tensors).
These encoders are shared across endpoints; they differ only in the semantic
role of the parents they read.
Second, fusion is \emph{endpoint-specific}: a learnable gate matrix
\(\in \mathbb{R}^{|\mathcal{V}_{\mathrm{target}}| \times B}\)
(with $B=4$ parent-type blocks, or $B=5$ when a triple block is enabled)
scales each type-level embedding separately for every target endpoint before
concatenation.
The gated embeddings are concatenated with the GNN target node embedding and
passed jointly to the classification head.
In \emph{typed+triple}, triple-path evidence is encoded by an additional
triple-path encoder and included as a fifth block, subject to its
own per-endpoint gate column in the endpoint-specific gate matrix.

\subsection{Non-binary parents}
Count- and continuous-valued direct parents are handled outside the binary
\(\varphi\)-coefficient framework. Values are mapped to comparable
pseudo-uniform scores via a train-only empirical CDF transform; count
variables receive deterministic, patient-ID-seeded jitter before
transformation. Non-linear dependence on continuous parents is represented
through a Legendre polynomial basis. These channels extend the binary \emph{patient activation}
in the 42-dimensional branch.

\subsection{Higher-order triple-path evidence}
To capture higher-order dependencies,
third-order moments are computed over graph-consistent
triples \(G \rightarrow P \rightarrow Y\).
The main coefficient \(a_{111}\) is estimated from a \(2 \times 2 \times 2\)
contingency table over three binary nodes on each DAG-consistent path
\(G \rightarrow P \rightarrow Y\) (grandparent, parent, and endpoint).
It generalizes the pairwise $\varphi$-coefficient, which uses only two
variables at a time.
The triple vector also retains the three pairwise marginals
(\(a_{11}^{GP}, a_{11}^{GY}, a_{11}^{PY}\)), prevalence terms, support,
patient activations, and an explicit patient evidence product
\(a_{111}\,\phi_G\,\phi_P\). 
Triple evidence uses the same slot-based encoder
interface as pairwise evidence.

The fusion modes connect the statistical branch to the GNN; they differ in
what evidence tensor is built and how it is combined with the deep
representation (\cref{tab:fusion_modes}).

\begin{table}[H]
  \centering
  \scriptsize
  \setlength{\tabcolsep}{4pt}
  \renewcommand{\arraystretch}{1.15}
  \caption{Fusion modes for the statistical evidence branch in Task~B.}
  \label{tab:fusion_modes}
  \begin{tabular}{@{}>{\raggedright\arraybackslash}p{0.14\textwidth}
                       >{\raggedright\arraybackslash}p{0.26\textwidth}
                       >{\raggedright\arraybackslash}p{0.28\textwidth}
                       >{\raggedright\arraybackslash}p{0.28\textwidth}@{}}
    \toprule
    \textbf{Mode} & \textbf{Evidence tensor} & \textbf{Encoder} & \textbf{Fusion with GNN} \\
    \midrule
    \texttt{linear}
      & Scalar HCR score
        \(s_{p,v}=\sum_r a_{11,rv}\,\phi_1(R_p)\) per endpoint
      & None (precomputed wide scores)
      & Added to endpoint logit;
        one learned scalar per endpoint
        (learned wide-branch weight) \\
    \texttt{nonlinear}
      & Pairwise 10D or 42D vectors per
        (patient, parent, endpoint) slot
      & Shared \texttt{HCRPairEncoder};
        mean-pool over parent slots
      & Encoder output added to logit;
        optional per-endpoint scale
        (learned per-endpoint scale) \\
    \texttt{typed\_concat}
      & 42D pairwise vectors, split into four
        type-specific tensors (\texttt{drug\_exposure},
        \texttt{mechanism},
        \texttt{adr\_or\_intermediate\_state},
        \texttt{patient\_context})
      & One pairwise evidence encoder per parent
        \emph{type} (shared across endpoints);
        endpoint-specific gate matrix
        $[|\mathcal{V}_{\text{target}}| \times 4]$
      & Gated type embeddings concatenated with
        GNN target embedding \emph{before}
        the classification head \\
    \texttt{triple}
      & 13-channel vectors per
        (patient, triple, endpoint) slot
        along paths \(G \rightarrow P \rightarrow Y\)
      & Shared pairwise evidence encoder on
        triple slots (replaces pairwise
        evidence in this mode)
      & Encoder output added to logit
        (no per-endpoint scalar) \\
    \emph{typed+triple}
      & Typed 42D pairwise blocks \emph{and}
        13-channel triple block
        (with triple-path evidence enabled)
      & Four typed encoders plus
        triple-path encoder; gates
        $[|\mathcal{V}_{\text{target}}| \times 5]$ (one column
        per block, including triple)
      & All gated blocks concatenated with GNN
        embedding before the head
        (same join as \texttt{typed\_concat}) \\
    \bottomrule
  \end{tabular}
\end{table}
\FloatBarrier

\section{Training and Evaluation}
\label{sec:taskB-train-eval}

\subsection{Training Objective}
\label{sec:taskB-train-obj}

Training iterates over patient graphs batched via PyG DataLoader.
Models were optimized using class-weighted binary cross-entropy with logits.
Endpoint weights were computed from the training-split prevalence, since clinical outcomes are rare and highly imbalanced.
Model selection uses ReduceLROnPlateau on validation performance.
Exploratory runs with focal loss were not retained in the final reported protocol.

\subsection{Evaluation Metrics}
\label{sec:taskB-metrics}

Evaluation is reported separately per endpoint and in aggregate, using AUROC, AP, and prevalence-normalized lift.
As the synthetic dataset generated for this project contains rare events and class imbalance, AP was chosen as a key performance metric.
It focuses on the positive class and reflects the trade-off between precision and recall.
The evaluation metrics are defined as follows.

\emph{AUROC} is the area under the receiver operating characteristic curve and is defined as:
\begin{equation}
\mathrm{AUROC}
=
\int_0^1 \mathrm{TPR}(\mathrm{FPR})\,d\mathrm{FPR},
\end{equation}
where $\mathrm{TPR}$ is the true positive rate and $\mathrm{FPR}$ is the false positive rate.

\emph{AP} is defined analogously like in Task~A (\cref{sec:taskA-ap}).

\emph{Lift} is the AP divided by the prevalence baseline, that is, the
factor by which the model improves over a random ranking, e.g.\ at a mean
prevalence of $5.5\%$ a lift of $2.3$ corresponds to an AP of $0.125$.

In addition to the traditional metrics described above, a
\emph{generator-based oracle AUROC} was computed for the node-classification
task. As the dataset is synthetic, the logistic generative model for
each endpoint is known: each patient is scored with the true parent effects
from the audited DAG and AUROC is evaluated on the resulting
probabilities. Model performance is reported relative to this reference via
\emph{\% of oracle AUROC}:

\begin{equation}
\frac{\mathrm{AUROC} - 0.5}{\mathrm{AUROC}_{\mathrm{Oracle}} - 0.5} \cdot 100\%,
\label{eq:pct-ceiling}
\end{equation}
%
where $\mathrm{AUROC}_{\mathrm{Oracle}} = 0.651$ is the macro-averaged oracle
AUROC over the same eight endpoints used in aggregate reporting
(\cref{tab:app_endpoints}). This metric expresses the share of ranking
signal recoverable under the known generator, relative to random performance
at $0.5$.

\subsection{Model and Training Configuration}
\label{sec:taskB-config}

The experimental configurations share a common training protocol.
All experiments share the same encoder structure and differ only in a small
number of hyperparameters: the number of layers $L$, the presence of a
residual connection and the statistical-evidence variant. Everything
else was held fixed. The shared configuration is listed in
\cref{tab:hparams-common}, and the parameters specific to individual
convolutions in \cref{tab:hparams-arch}.

% =====================================================================
%  HYPERPARAMETERS
% =====================================================================
\begin{table}[H]
  \centering
  \caption{Task~B: shared model and training configuration.
    Values were fixed across the experiments reported in this chapter.}
  \label{tab:hparams-common}
  \begin{tabular}{lll}
    \toprule
    Group & Hyperparameter & Value \\
    \midrule
    \multirow{5}{*}{Encoder}
      & Hidden units per convolutional layer & 64 \\
      & Number of layers $L$                 & 1, 2, 3, 4, 5, 6 \\
      & Aggregation across relations         & sum \\
      & Activation                           & ReLU \\
      & Batch normalisation                  & yes \\
    \midrule
    \multirow{3}{*}{Pair encoder}
      & Hidden units                         & 32 \\
      & Output embedding                     & 8 \\
      & Dropout                              & 0.0 \\
    \midrule
    \multirow{2}{*}{Regularisation}
      & Dropout                              & 0.2 \\
      & Weight decay                         & $5\times10^{-4}$ \\
    \midrule
    \multirow{3}{*}{Classification head}
      & Number of layers                     & 2 \\
      & Hidden units                         & 64 (equal to encoder width) \\
      & Output units                         & 1 per endpoint \\
    \midrule
    \multirow{8}{*}{Optimisation}
      & Optimiser                            & Adam \\
      & Learning rate                        & $1\times10^{-3}$ \\
      & Batch size                           & 128 \\
      & Epochs (maximum)                     & 200 \\
      & Loss                                 & weighted BCE with logits \\
      & Positive weight cap                  & 30 \\
      & Gradient clipping                    & disabled \\
      & Seed                                 & 42 \\
    \midrule
    \multirow{3}{*}{Model selection}
      & Selection metric                     & validation AP \\
      & Early-stopping patience              & 20 epochs \\
      & Minimum improvement                  & 0.001 \\
    \bottomrule
  \end{tabular}
\end{table}

    \begin{table}[H]
      \centering
      \small
      \caption{Task~B: architecture-specific configuration.}
      \label{tab:hparams-arch}
      \begin{tabularx}{\textwidth}{l >{\raggedright\arraybackslash}X
                                     >{\raggedright\arraybackslash}X
                                     >{\raggedright\arraybackslash}X}
        \toprule
        Hyperparameter & R-GCN & GraphSAGE & TransformerConv \\
        \midrule
        Attention heads              & --- & --- & 4 \\
        Units per head               & --- & --- & 16 \\
        Output width per relation    & 64 & 64 & $4\times16=64$ \\
        Basis decomposition          & 8 bases & --- & --- \\
        Uses edge attributes         & relation id only & none (stored, unused) & full \texttt{edge\_attr} \\
        Weight matrices per relation & shared through bases & 2 & 3 (query, key, value) \\
        \bottomrule
      \end{tabularx}
    \end{table}

All three convolutions are wrapped in \texttt{HeteroConv}, which instantiates a separate
    convolution for each of the 19 relation types in the graph. The
    parameter counts therefore scale with the number of relations. In
    attention-based layers the hidden width is split across heads, so the
    output width stays at 64 regardless of the number of heads.

The entries in \cref{tab:hparams-common,tab:hparams-arch} were held fixed so that
the reported comparisons isolate depth, residual connections, and the
statistical-evidence variant.
Adam with learning rate \(10^{-3}\) follows the standard setting of
Kingma and Ba~\cite{kingma2015adam}; ReLU, batch normalisation, and the
64-dimensional width are the same conventions as in the graph encoder, so
that R-GCN, GraphSAGE, and TransformerConv differ only in the
message-passing operator.
\texttt{HeteroConv} uses sum aggregation across relation types after
preliminary checks in which mean and max aggregation did not change the
conclusions.
The pair-encoder widths are those selected in \cref{subsec:hcr_application}.
Batch size 128 was used as a stable setting for patient-graph batches; the
remaining optimisation details (including the choice of Adam) were not
ablated.

\section{Experimental Results}
\label{sec:taskB-results}

\subsection{Architecture Comparison}
\label{sec:taskB-depth}

The first set of experiments established how much the choice of
GNN architecture matters for this task. Three message-passing layers were
compared: R-GCN, GraphSAGE and TransformerConv. All three are defined for
homogeneous graphs, so each was wrapped in \texttt{HeteroConv}, which
instantiates a separate embedding for every relation type and aggregates
the results at the target node and allows to preserve the heterogenous structure.

The results for those experiments are reported in \cref{tab:arch-nores-update}.
All of metrics reported are described in \cref{subsubsec:oversmoothing}.
The epoch is selected on the validation set and
test metrics are reported from that epoch.

    \begin{table}[H]
      \centering
      \small
      \caption{Task~B: message-passing architecture comparison without residual connections.}
      \label{tab:arch-nores-update}
      {\sisetup{table-format=1.3}
      \resizebox{\textwidth}{!}{%
      \begin{tabular}{l c SSS S[table-format=2.1] SSS S[table-format=2.1]}
        \toprule
        & & \multicolumn{4}{c}{Validation} & \multicolumn{4}{c}{Test} \\
        \cmidrule(lr){3-6}\cmidrule(lr){7-10}
        Architecture & $L$ & {AUROC} & {AP} & {Lift} & {\% oracle} & {AUROC} & {AP} & {Lift} & {\% oracle} \\
        \midrule
        R-GCN            & 1 & 0.623 & 0.124 & 2.28 & 81.6 & 0.624 & 0.120 & 2.21 & 82.0 \\
        GraphSAGE        & 1 & 0.622 & 0.124 & 2.28 & 80.9 & 0.619 & 0.120 & 2.21 & 78.9 \\
        TransformerConv  & 1 & 0.622 & 0.125 & 2.30 & 81.2 & 0.620 & 0.120 & 2.20 & 79.6 \\
        \addlinespace[2pt]
        R-GCN            & 3 & 0.542 & 0.076 & 1.40 & 28.0 & 0.542 & 0.068 & 1.26 & 27.7 \\
        GraphSAGE        & 3 & 0.549 & 0.079 & 1.45 & 32.2 & 0.549 & 0.072 & 1.32 & 32.8 \\
        TransformerConv  & 3 & 0.550 & 0.083 & 1.53 & 33.2 & 0.552 & 0.073 & 1.35 & 34.5 \\
        \addlinespace[2pt]
        R-GCN            & 6 & 0.503 & 0.055 & 1.00 &  2.2 & 0.500 & 0.054 & 1.00 &  0.0 \\
        GraphSAGE        & 6 & 0.530 & 0.065 & 1.20 & 20.0 & 0.510 & 0.061 & 1.12 &  6.6 \\
        TransformerConv  & 6 & 0.527 & 0.065 & 1.20 & 18.0 & 0.500 & 0.054 & 1.00 &  0.0 \\
        \bottomrule
      \end{tabular}}}
    \end{table}

The experiments were conducted with multiple number of layers as 
patient path is up to 6 degrees and it is assumed that every point
can have an influence in casual path. For each experiment all number 
of layers were calculated, 
but for reporting reasons in \cref{tab:arch-nores-update} there are included layers 
1,3 and 6 to show the trend with the increased number of layers. The reporting 
convention is that unless 
the results from other layers are significant for conclusion. 

% The range of depths follows the assumption that 
% %the longest path from a drug exposure to an endpoint in
% %the audited DAG spans up to six hops, and 
% every node along such a path may
% contribute to the causal chain. Models with one to six layers were therefore
% trained for every configuration. For readability only $L \in \{1, 3, 6\}$ are
% listed in \cref{tab:arch-nores-update}, which is sufficient to show the trend.
% Intermediate depths are reported only where they are relevant to a specific
% conclusion.

Without residual connections the three architectures diverge sharply at
L=6, collapse to constant predictions (AUROC=0.50).

\subsection{Effect of Residual Connections}
\label{sec:taskB-residual}

That's why the next step was application of residual connections
    to see the capabilities of models with higher number of layers. The results are presented 
    in \cref{tab:arch-res}.
    
    \begin{table}[H]
      \centering
      \small
      \caption{Task~B: message-passing architecture comparison with residual connections.}
      \label{tab:arch-res}
      {\sisetup{table-format=1.3}
      \resizebox{\textwidth}{!}{%
      \begin{tabular}{l c SSS S[table-format=2.1] SSS S[table-format=2.1]}
        \toprule
        & & \multicolumn{4}{c}{Validation} & \multicolumn{4}{c}{Test} \\
        \cmidrule(lr){3-6}\cmidrule(lr){7-10}
        Architecture & $L$ & {AUROC} & {AP} & {Lift} & {\% oracle} & {AUROC} & {AP} & {Lift} & {\% oracle} \\
        \midrule
        R-GCN            & 1 & 0.625 & 0.125 & 2.30 & 82.7 & 0.625 & 0.121 & 2.22 & 83.2 \\
        GraphSAGE        & 1 & 0.621 & 0.124 & 2.27 & 80.6 & 0.619 & 0.120 & 2.21 & 78.7 \\
        TransformerConv  & 1 & 0.623 & 0.126 & 2.31 & 81.3 & 0.619 & 0.120 & 2.21 & 79.0 \\
        \addlinespace[2pt]
        R-GCN            & 3 & 0.630 & 0.133 & 2.44 & 86.0 & 0.627 & 0.125 & 2.29 & 84.2 \\
        GraphSAGE        & 3 & 0.623 & 0.132 & 2.42 & 81.4 & 0.616 & 0.124 & 2.27 & 77.0 \\
        TransformerConv  & 3 & 0.628 & 0.135 & 2.47 & 85.2 & 0.620 & 0.123 & 2.27 & 79.7 \\
        \addlinespace[2pt]
        R-GCN            & 6 & 0.629 & 0.135 & 2.48 & 85.6 & 0.626 & 0.123 & 2.27 & 83.9 \\
        GraphSAGE        & 6 & 0.628 & 0.134 & 2.46 & 84.9 & 0.624 & 0.127 & 2.33 & 82.3 \\
        TransformerConv  & 6 & 0.627 & 0.134 & 2.46 & 84.2 & 0.620 & 0.122 & 2.24 & 79.4 \\
        \bottomrule
      \end{tabular}}}
    \end{table}

    With residual connections all three architectures converge to comparable
    performance (\cref{tab:arch-res}): 77--86\% of the oracle AUROC
    (\cref{eq:pct-ceiling}) on the test
    set, with differences of about 0.01 AUROC. The more informative factors for model assessment are AP and 
    Lift, which show that the model performs nearly 2.3 times better than a random baseline. 
    Since the architectures achieve comparable performance, the choice of architecture
    for 
    the remaining experiments 
    was therefore made on structural grounds. TransformerConv is the only one of the three
    backbones benchmarked here that consumes the full \texttt{edge\_attr} vector in message
    passing (see \cref{tab:edge_attr_by_conv}).
    %signed effect
    %(\texttt{effect\_size}$\times$\texttt{effect\_sign}), audit scalars
    %(\texttt{activation\_frequency}, \texttt{mechanistic\_confidence},
    %\texttt{evidence\_weight}), and the one-hot \texttt{edge\_type} block.
   %GraphSAGE ignores edge features; R-GCN conditions on relation identity
  %hrough separate relation modules rather than on \texttt{edge\_attr}.
    Since edge-level audit evidence is important for graph representation, 
    TransformerConv was
    selected for all subsequent experiments with the statistical evidence branch.

\subsection{Contribution of Statistical Evidence}
\label{sec:taskB-hcr-results}

The statistical evidence branch was introduced to supply the model with empirical dependence
signals about grandparent/parent--endpoint pairs that the message-passing path has to infer on
its own (\cref{subsec:hcr_application}). The configurations examined here
vary along three axes, each addressing a separate question. The first is the
form of the wide branch: a single learned coefficient per endpoint applied
to a precomputed dependence score, against a small nonlinear encoder applied
to each parent separately. The second is the scope of the evidence: binary
parents only, or all node types, in which case continuous and count
variables are expanded in a Legendre basis. The third is the order of the
statistic: pairwise coefficients between a parent and an endpoint, or
third-order moments over grandparent--parent--endpoint paths, which capture
interactions that pairwise coefficients cannot represent. The relevant 
Abbreviations used in the result tables are given in \cref{tab:hcr-legend};
the corresponding scores are reported in \cref{tab:hcr-nores}.

\begin{table}[H]
  \centering
  \small
  \caption{Notation used for the statistical-evidence variants in the result tables.}
  \label{tab:hcr-legend}
  \begin{tabularx}{\textwidth}{l l >{\raggedright\arraybackslash}X}
    \toprule
    Label & Join point & Description \\
    \midrule
    ---
      & {--}
      & No statistical evidence branch \\
    \addlinespace[3pt]
    linear
      & sum at logit
      & Precomputed dependence score, all parent types; one learned scalar
        per endpoint, no encoder \\
    \addlinespace[3pt]
    nonlin.\ all 42D
      & sum at logit
      & 42 evidence channels, all parent types; shared pair encoder \\
    \addlinespace[3pt]
    nonlin.\ bin.\ 10D
      & sum at logit
      & 10 evidence channels, binary parents only; shared pair encoder \\
    \addlinespace[3pt]
    nonlin.\ bin.\ 42D
      & sum at logit
      & 42 evidence channels, binary parents only; shared pair encoder \\
    \addlinespace[3pt]
    triple
      & sum at logit
      & 13 channels of third-order moments over
        grandparent--parent--endpoint paths, replacing the pairwise
        evidence \\
    \addlinespace[3pt]
    typed+triple
      & concatenation
      & 42 pairwise channels split across four non-shared encoders, one per
        parent type, each with a learned per-endpoint gate, plus a
        third-order block with its own per-endpoint gate column \\
    \bottomrule
  \end{tabularx}
\end{table}

\begin{table}[H]
  \centering
  \small
  \setlength{\tabcolsep}{4pt}
  \caption{Task~B: statistical-evidence variants without residual connections
    (TransformerConv backbone).}
  \label{tab:hcr-nores}
  \begin{tabular}{clcccccccc}
    \toprule
    & & \multicolumn{4}{c}{Validation} & \multicolumn{4}{c}{Test} \\
    \cmidrule(lr){3-6}\cmidrule(lr){7-10}
    $L$ & Evidence variant & AUROC & AP & Lift & \% oracle & AUROC & AP & Lift & \% oracle \\
    \midrule
    1 & ---              & 0.622 & 0.1251 & 2.30 & 81.19 & 0.620 & 0.1197 & 2.20 & 79.61 \\
    1 & linear 42D       & 0.622 & 0.1247 & 2.29 & 81.10 & 0.617 & 0.1191 & 2.19 & 77.83 \\
    1 & nonlin. all 42D  & 0.622 & 0.1243 & 2.28 & 81.17 & 0.618 & 0.1192 & 2.19 & 78.57 \\
    1 & typed+triple     & 0.634 & \textbf{0.1348} & 2.48 & \textbf{89.05} & 0.626 & 0.1241 & 2.28 & \textbf{83.46} \\
    \addlinespace[4pt]
    2 & nonlin. all 42D  & 0.625 & 0.1242 & 2.28 & 82.77 & 0.615 & 0.1180 & 2.17 & 76.15 \\
    2 & typed+triple     & 0.628 & 0.1328 & 2.44 & 84.91 & 0.621 & 0.1235 & 2.27 & 80.16 \\
    \addlinespace[4pt]
    3 & ---              & 0.550 & 0.0833 & 1.53 & 33.19 & 0.552 & 0.0734 & 1.35 & 34.49 \\
    3 & linear 42D       & 0.604 & 0.1173 & 2.15 & 69.28 & 0.612 & 0.1031 & 1.90 & 74.51 \\
    3 & triple           & 0.602 & 0.1131 & 2.08 & 67.92 & 0.590 & 0.1021 & 1.88 & 59.81 \\
    3 & nonlin. all 42D  & 0.621 & 0.1310 & 2.41 & 80.13 & 0.629 & 0.1197 & 2.20 & 85.63 \\
    3 & typed+triple     & 0.625 & \textbf{0.1362} & 2.50 & \textbf{82.90} & 0.620 & \textbf{0.1252} & 2.30 & 79.54 \\
    \addlinespace[4pt]
    4 & nonlin. all 42D  & 0.616 & 0.1214 & 2.23 & 76.90 & 0.609 & 0.1168 & 2.15 & 72.30 \\
    4 & typed+triple     & 0.631 & 0.1350 & 2.48 & 86.65 & 0.614 & 0.1237 & 2.27 & 75.53 \\
    \addlinespace[4pt]
    5 & nonlin. all 42D  & 0.617 & 0.1183 & 2.17 & 77.67 & 0.611 & 0.1132 & 2.08 & 73.37 \\
    5 & typed+triple     & 0.627 & 0.1320 & 2.42 & 84.31 & 0.617 & 0.1225 & 2.25 & 77.58 \\
    \addlinespace[4pt]
    6 & ---              & 0.527 & 0.0653 & 1.20 & 17.96 & 0.500 & 0.0544 & 1.00 &  0.00 \\
    6 & linear 42D       & 0.616 & 0.1206 & 2.21 & 76.91 & 0.603 & 0.0955 & 1.76 & 68.22 \\
    6 & triple           & 0.607 & 0.1154 & 2.12 & 70.71 & 0.579 & 0.0967 & 1.78 & 52.56 \\
    6 & nonlin. all 42D  & 0.569 & 0.0955 & 1.75 & 45.53 & 0.500 & 0.0544 & 1.00 &  0.00 \\
    6 & typed+triple     & 0.630 & 0.1326 & 2.43 & 86.53 & 0.618 & 0.1218 & 2.24 & 78.35 \\
    \bottomrule
  \end{tabular}
\end{table}

Without the statistical evidence branch the model degrades sharply with depth and predicts a constant at
$L=6$. Every statistical-evidence variant recovers a substantial part of the lost performance,
and the gain grows with depth: $+1.8$, $+42.5$ and $+74.5$ percentage points
of the oracle AUROC (\cref{eq:pct-ceiling}) at $L=1$, $3$ and $6$ respectively.
The concatenated
variant is the only one that remains flat across all six depths.

The reason for improved results with the statistical evidence branch is that it reads the
dependence statistics of the direct parents and grandparents from a precomputed tensor and
never passes through a graph convolution, so its contribution is identical at
$L=1$ and at $L=6$. 

% =====================================================================
%  HCR variants on TransformerConv, synthetic dataset.
% =====================================================================

% \begin{table}[htbp]
%   \centering
%   \small
%   \setlength{\tabcolsep}{4pt}
%   \caption{HCR variants on TransformerConv \emph{without} residual
%     connections.}
%   \label{tab:hcr-nores}
%   \begin{tabular}{clcccccccc}
%     \toprule
%     & & \multicolumn{4}{c}{Validation} & \multicolumn{4}{c}{Test} \\
%     \cmidrule(lr){3-6}\cmidrule(lr){7-10}
%     $L$ & HCR & AUC & AUPRC & Lift & \% oracle & AUC & AUPRC & Lift & \% oracle \\
%     \midrule
%     1 & ---            & 0.622 & 0.1251 & 2.30 & 81.19 & 0.620 & 0.1197 & 2.20 & 79.61 \\
%     1 & linear 42D     & 0.622 & 0.1247 & 2.29 & 81.10 & 0.617 & 0.1191 & 2.19 & 77.83 \\
%     1 & typed+triple   & \textbf{0.630} & 0.1331 & 2.44 & \textbf{86.30} & \textbf{0.623} & 0.1229 & 2.26 & \textbf{81.40} \\
%     \addlinespace[4pt]
%     2 & typed+triple   & 0.628 & 0.1314 & 2.41 & 84.67 & 0.621 & 0.1228 & 2.26 & 80.13 \\
%     \addlinespace[4pt]
%     3 & ---            & 0.550 & 0.0833 & 1.53 & 33.19 & 0.552 & 0.0734 & 1.35 & 34.49 \\
%     3 & linear 42D     & 0.604 & 0.1173 & 2.15 & 69.28 & 0.612 & 0.1031 & 1.90 & 74.51 \\
%     3 & triple         & 0.602 & 0.1131 & 2.08 & 67.92 & 0.590 & 0.1021 & 1.88 & 59.81 \\
%     3 & typed+triple   & 0.618 & 0.1265 & 2.32 & 78.41 & 0.616 & 0.1214 & 2.23 & 77.02 \\
%     \addlinespace[4pt]
%     4 & typed+triple   & \textbf{0.630} & \textbf{0.1369} & \textbf{2.51} & 86.09 & 0.616 & \textbf{0.1239} & \textbf{2.28} & 76.97 \\
%     \addlinespace[4pt]
%     5 & typed+triple   & 0.627 & 0.1308 & 2.40 & 84.25 & 0.617 & 0.1222 & 2.25 & 77.61 \\
%     \addlinespace[4pt]
%     6 & ---            & 0.527 & 0.0653 & 1.20 & 17.96 & 0.500 & 0.0544 & 1.00 &  0.00 \\
%     6 & linear 42D     & 0.616 & 0.1206 & 2.21 & 76.91 & 0.603 & 0.0955 & 1.76 & 68.22 \\
%     6 & triple         & 0.607 & 0.1154 & 2.12 & 70.71 & 0.579 & 0.0967 & 1.78 & 52.56 \\
%     6 & typed+triple   & 0.622 & 0.1222 & 2.24 & 80.91 & 0.612 & 0.1179 & 2.17 & 74.53 \\
%     \bottomrule
%   \end{tabular}
% \end{table}

    Residual connections address the same failure from the opposite side: they
    preserve the signal inside the propagation channel instead of bypassing it.
    Both residual connections and the statistical-evidence variants are effective on this benchmark, but they are not complementary in a consistent way.
    With residual connections enabled, the evaluated statistical-evidence variants did not provide a consistent additional improvement over the corresponding graph-only residual models on the synthetic benchmark
    (\cref{tab:hcr-res}).

    \begin{table}[H]
      \centering
      \small
      \setlength{\tabcolsep}{4pt}
      \caption{Task~B: statistical-evidence variants with residual connections
        (TransformerConv backbone).}
      \label{tab:hcr-res}
      \begin{tabular}{clcccccccc}
        \toprule
        & & \multicolumn{4}{c}{Validation} & \multicolumn{4}{c}{Test} \\
        \cmidrule(lr){3-6}\cmidrule(lr){7-10}
        $L$ & Evidence variant & AUROC & AP & Lift & \% oracle & AUROC & AP & Lift & \% oracle \\
        \midrule
        1 & ---            & 0.623 & 0.1255 & 2.31 & 81.30 & 0.619 & 0.1202 & 2.21 & 79.02 \\
        1 & nonlin. bin.   & 0.622 & 0.1238 & 2.27 & 81.12 & 0.617 & 0.1176 & 2.16 & 77.82 \\
        1 & nonlin. 42D    & 0.621 & 0.1250 & 2.30 & 80.56 & 0.616 & 0.1186 & 2.18 & 76.95 \\
        1 & typed+triple   & \textbf{0.631} & 0.1335 & 2.45 & \textbf{86.94} & \textbf{0.621} & 0.1240 & 2.28 & \textbf{80.30} \\
        \addlinespace[4pt]
        2 & typed+triple   & 0.631 & 0.1322 & 2.43 & 86.77 & 0.619 & 0.1230 & 2.26 & 79.05 \\
        \addlinespace[4pt]
        3 & ---            & 0.628 & 0.1346 & 2.47 & 85.23 & 0.620 & 0.1233 & 2.27 & 79.66 \\
        3 & nonlin. bin.   & 0.627 & 0.1339 & 2.46 & 84.06 & 0.617 & 0.1242 & 2.28 & 77.48 \\
        3 & nonlin. 42D    & 0.630 & 0.1343 & 2.47 & 85.95 & 0.619 & \textbf{0.1247} & 2.29 & 78.80 \\
        3 & typed+triple   & 0.629 & 0.1350 & 2.48 & 85.64 & 0.618 & 0.1246 & 2.29 & 78.05 \\
        \addlinespace[4pt]
        4 & typed+triple   & 0.627 & \textbf{0.1367} & \textbf{2.51} & 84.10 & 0.621 & \textbf{0.1256} & \textbf{2.31} & 80.08 \\
        \addlinespace[4pt]
        5 & typed+triple   & 0.626 & 0.1361 & 2.50 & 83.88 & \textbf{0.623} & \textbf{0.1256} & \textbf{2.31} & \textbf{81.54} \\
        \addlinespace[4pt]
        6 & ---            & 0.627 & 0.1337 & 2.46 & 84.19 & 0.620 & 0.1219 & 2.24 & 79.37 \\
        6 & nonlin. bin.   & 0.614 & 0.1177 & 2.16 & 75.90 & 0.605 & 0.1082 & 1.99 & 69.84 \\
        6 & nonlin. 42D    & 0.624 & 0.1301 & 2.39 & 82.21 & 0.618 & 0.1214 & 2.23 & 78.23 \\
        6 & typed+triple   & 0.629 & 0.1347 & 2.47 & 85.62 & 0.620 & 0.1244 & 2.29 & 79.34 \\
        \bottomrule
      \end{tabular}
    \end{table}
    
%triple najlepsze i dodac do tego te macierz sciezek 
%dodac wyniki per endpoint

For understanding how the statistical evidence branch contributes to individual endpoint prediction, the learned
gate weights are shown in \cref{fig:influence}. Each cell gives the scalar by
which the model scales the evidence block of one parent type before it enters
the classification head; the gates were initialised at $1.0$, so values below
unity indicate reduced reliance on that channel.

\begin{figure}[H]
  \centering
  \includegraphics[width=\textwidth]{hcr_gate_matrix.pdf}
  \caption{Statistical-evidence gate matrix: influence of parent node type per endpoint.}
  \label{fig:influence}
\end{figure}

The gate ordering aligns with the number of endpoints that
have at least one direct parent of each type in the graph,
suggesting that the model allocates more weight to parent types
that are structurally more informative.

\subsection{Final Evaluation}
\label{sec:taskB-final}

In summary, models with a statistical evidence branch reached results
comparable to models with residual connections applied. 
The gate matrix indicates that the evidence branch is used
unevenly across endpoints and parent types, in line with differences in
parent availability and predictive utility on the synthetic task. 
In general the value of the statistical evidence branch therefore does not lie in raising the peak, but
in bringing stability to learning across all depths and additional information useful for clinical
event prediction.
Across the six depths evaluated, \emph{typed+triple} is the only variant that
never collapses, retaining between $82\%$ and $89\%$ of the oracle AUROC, while
both the baseline and \emph{nonlin.\ all 42D} fall below $80.00\%$ at $L=6$. 
The configuration \emph{typed+triple} is mostly stable on depth
and additionally brings the most informative value.
The highest validation score in \cref{tab:hcr-nores} is obtained by
\emph{typed+triple} at $L=1$ and $L=4$ ($89.05\%$ and $86.65\%$ of the oracle AUROC). 
Taken together, the \emph{typed+triple} statistical-evidence variant is selected as the preferred configuration.

In \cref{tab:per-endpoint} there are presented results per endpoint.

    \begin{table}[H]
      \centering
      \small
      \setlength{\tabcolsep}{4pt}
      \caption{Task~B: per-endpoint results for the selected model
        (\emph{typed+triple} at $L=4$ without residual connections).}
      \label{tab:per-endpoint}
      \begin{tabular}{lccc ccc ccc}
        \toprule
        & & \multicolumn{2}{c}{Oracle} & \multicolumn{3}{c}{Validation} & \multicolumn{3}{c}{Test} \\
        \cmidrule(lr){3-4}\cmidrule(lr){5-7}\cmidrule(lr){8-10}
        Endpoint & Prev. & AUROC & AP & AUROC & AP & \%~oracle & AUROC & AP & \%~oracle \\
        \midrule
        AKI                & 0.131 & 0.811 & 0.454 & 0.794 & 0.406 & 94.3  & 0.785 & 0.421 & 91.6 \\
        Depression         & 0.085 & 0.652 & 0.166 & 0.655 & 0.165 & 101.6 & 0.639 & 0.170 & 91.0 \\
        Hyponatremia       & 0.079 & 0.638 & 0.168 & 0.613 & 0.132 & 81.5  & 0.607 & 0.142 & 77.5 \\
        Hyperkalemia       & 0.049 & 0.590 & 0.083 & 0.595 & 0.080 & 105.7 & 0.550 & 0.074 & 55.2 \\
        Delirium           & 0.037 & 0.668 & 0.100 & 0.649 & 0.093 & 89.0  & 0.687 & 0.098 & 93.0 \\
        GI bleeding        & 0.021 & 0.570 & 0.080 & 0.570 & 0.080 & 99.0  & 0.528 & 0.029 & 39.6 \\
        DILI               & 0.018 & 0.653 & 0.086 & 0.563 & 0.067 & 41.3  & 0.579 & 0.036 & 51.6 \\
        QT arrhythmia      & 0.016 & 0.622 & 0.075 & 0.591 & 0.050 & 74.1  & 0.536 & 0.019 & 29.4 \\
        \midrule
        \textbf{Macro}     & 0.055 & 0.651 & 0.151 & 0.629 & 0.134 & 85.4  & 0.614 & 0.124 & 75.5 \\
        \midrule
        \multicolumn{10}{l}{\emph{Trained but excluded from the macro average}} \\
        Lactic acidosis    & 0.0057 & 0.523 & 0.011 & 0.560 & 0.009 & {--} & 0.549 & 0.008 & {--} \\
        Rhabdomyolysis     & 0.0037 & 0.530 & 0.011 & 0.536 & 0.005 & {--} & 0.724 & 0.033 & {--} \\
        Serotonin syndrome & 0.0027 & 0.554 & 0.008 & 0.646 & 0.010 & {--} & 0.642 & 0.008 & {--} \\
        \bottomrule
      \end{tabular}
    \end{table}
